%_____________________________________________________________________________________________ 
% Chapter 1: Introduction
%_____________________________________________________________________________________________ 
\chapter{Introduction}
\label{sec:introduction}

Anyone who has tried to turn a BFT consensus paper into running code knows the frustration:
the protocol itself may be clean, but everything around it---storage, configuration changes,
application hooks, and keeping the system alive when nobody is sending requests---gets
hand-waved away.  HotStuff~\cite{hotstuff2019} is a good example.  On paper it is one of
the tidiest BFT protocols available: $O(n)$ messages per decision, optimistic
responsiveness, and a 3-chain commit rule that is easy to reason about.  Meta trusted it
enough to use it as the consensus engine for Diem.  But what about the publicly available
prototypes?  They crash and lose everything.  They assume that the validator set never
changes.  They hard-code a trivial state machine, and if clients stop sending requests for a
while, the protocol stalls.

We are not the first to notice.  A recent Stanford CS244B course
report~\cite{stanford_hotstuff2024} catalogued exactly these four pain points after the
authors tried to build a working HotStuff system from scratch.

\begin{enumerate}
  \item \textbf{No persistence.}  State lives in RAM.  Kill a replica and it forgets
    everything, including but not limited to committed blocks, votes, and quorum.
  \item \textbf{Static membership.}  The validator set is baked in at startup.  If a fifth
    node needs to be added, we would need to restart the whole cluster.
  \item \textbf{Hard-coded state machine.}  The ``application'' is usually an append-only
    log with no hooks for real business logic.
  \item \textbf{Liveness stalls under low load.}  The 3-chain rule needs three consecutive
    justified blocks.  No client requests means no new blocks, which means nothing ever
    commits.
\end{enumerate}

We started this project because we wanted a HotStuff implementation that addresses all four
issues in one coherent system.  The result is \textit{PersistHotStuff}, whose design
proposal we will present here, backed by implementation in Rust~\cite{rust-implementation}.

\section{Contributions}

Our contributions, then, are:

\begin{itemize}
  \item A \textbf{Persistence layer} made up of Write Ahead Logs (WAL) + snapshots adapted
    from Raft~\cite{raft2014} to the Byzantine setting.
  \item A \textbf{feature for dynamic membership} via consensus-driven reconfiguration
    commands.
  \item A \textbf{pluggable state machine interface} (trait-based, with gRPC for
    out-of-process applications).
  \item A \textbf{pacemaker design} that uses dummy proposals to keep the 3-chain growing
    during idle periods.
  \item An \textbf{open-source Rust codebase} with Ed25519 authentication, a simulation
    framework, and preliminary test results.
\end{itemize}

\section{Report Organisation}

The rest of this report is structured as follows.
Chapter~\ref{sec:background} provides background on the system model and HotStuff mechanics.
Chapter~\ref{sec:related_work} surveys related work.
Chapter~\ref{sec:system_design} presents the system design in detail.
Chapter~\ref{sec:evaluation} reports evaluation results---simulation, stress tests,
formal verification, and performance benchmarks.
Chapter~\ref{sec:discussion} discusses design trade-offs and safety/liveness arguments.
Chapter~\ref{sec:conclusion} concludes.
